%Plantilla para escritura de seminario/tesis
%Universidad Catolica del Norte, Chile
%Escuela de Ciencias Empresariales
%Modificada de 
%Iván Cornejo www.cornejog.com/memoristas
%2021-03-01
\documentclass[12pt, oneside]{report}

\usepackage[natbibapa]{apacite}
\usepackage[spanish]{babel}
\usepackage{mathptmx}
\usepackage[shortlabels]{enumitem}
%\usepackage{lmodern}
\usepackage[most]{tcolorbox}
\usepackage{fontawesome5}
\usepackage{smartdiagram}
\usepackage{rotating}
\usepackage{multirow}
\usepackage{amsmath}
\DeclareMathOperator{\re}{Re}
\usepackage{pgfplots}
\pgfplotsset{compat=1.16}
\usepackage{caption}
\usepackage{subcaption}
\usepackage[hidelinks]{hyperref}

 

\usepackage{titlesec}

\titleformat{\chapter}{\normalfont\Large\bf\MakeUppercase}{\thechapter.}{10pt}{\Large\bf\MakeUppercase}
\titlespacing*{\chapter}{0pt}{0pt}{40pt}

\usepackage[letterpaper,left=3.5cm,right=3cm,top=3cm,bottom=3cm]{geometry} %control de tamaño de papel y márgenes
\usepackage{graphicx} %para poder incluir figuras 
\usepackage{siunitx} %para notación científica
\usepackage{float} %complemento para figura
\usepackage{import} %para poder dividir el texto en varios archivos
\usepackage{natbib} %para la bibliografía
\usepackage{setspace} \onehalfspacing %Control del espacio entre lineas

%Hoja de la portada
\title{
\vspace{-3.0cm}
\includegraphics[width=2cm,height=2cm]{figuras/Escudo-UCN-Full-Color.png} \hspace{6.92cm}\includegraphics[width=6cm,height=2.cm]{figuras/logo_eciem.png}\\
\vspace{0.5cm}
\large \textbf{UNIVERSIDAD CATÓLICA DEL NORTE}

\normalsize\textbf{ESCUELA DE CIENCIAS EMPRESARIALES}

\small\textbf{Ingenier\'ia en Informaci\'on y Control de Gesti\'on}

\vspace{3.cm}
%{\Large Seminario Su perior en Sistemas}\\
{\LARGE \textbf{\uppercase{``Titulo del seminario ''}}}

\vspace{0.5cm}
{\footnotesize \begin{flushright} \textit{Preparaci\'on de Actividad de Título, parte de la actividad de Titlulaci\'on } \\ \textit{para optar al Título profesional de Ingeniero} 

\textit{en Información y Control de Gestión.} \end{flushright}}
\vspace{0.5cm}
{\footnotesize \begin{flushright} \textit{\textbf{Profesor Guía: Nombre profesor guía}} \end{flushright}}
\vspace{1.5cm}%CAMBIAR EL ESPACIO DE ACUERDO AL TITULO
{\normalsize \centering
    \textbf{autor 1} 
    
    autor1@alumnos.ucn.cl
    
    \textbf{autor2}
    
    autor2@alumnos.ucn.cl
    
    \textbf{autor3} 
    
    autor3@alumnos.ucn.cl
    
    }
    
\vspace{0.5cm}
\date{\textit{\centering \footnotesize Coquimbo, \footnotesize{\today}}}
}

\author{} %dejar vacío
\date{} %dejar vacío
%----------

%Este bloque es para adaptar los nombres a español

%----------

%Control del encabezado y pie de página
\usepackage{fancyhdr}
\pagestyle{fancy}
\setlength{\headheight}{15pt}
\renewcommand{\headrulewidth}{0pt}
\fancyhf{}
\rhead{}
\lhead{}
\fancyfoot[R]{\thepage}
\begin{document}
\pagenumbering{roman} %comienza con número romano para las páginas
\maketitle %agrega la portada al documento

\renewcommand{\bibname}{Bibliografía}
\renewcommand{\figurename}{Figura}
\renewcommand{\tablename}{Tabla}
\renewcommand{\listfigurename}{Índice de Figuras}
\renewcommand{\listtablename}{Índice de Tablas}
\renewcommand{\contentsname}{Índice}
\renewcommand{\chaptername}{Capítulo}
%se incluye lo contenido en resumen.tex
%\include{resumen}
%se incluye lo contenido en abstract.tex
%\include{abstract}
%se incluye lo contenido en agradecimientos.tex
%\include{agradecimientos}
%Estas 3 listas se crearán automáticamente, según los capítulos y secciones del documento 
\tableofcontents    %crea la lista de contenidos
\listoftables       %crea la lista de tablas
\listoffigures      %crea la lista de figuras
%incluye lo contenido en nomenclatura.tex
% \include{nomenclatura}
%incluye lo contenido en C1problema.tex
\include{C1problema}
%incluye lo contenido en C3objetivos.tex
\include{C3objetivos}
%incluye lo contenido en C2marcoteorico.tex
\include{C2marcoteorico}
%incluye lo contenido en C4conclusiones.tex
\include{C4medologia}
%incluye lo contenido en C5cartagantt.tex para preparacion de seminario
%include{C5CartaGantt}
\include{C5resultados} %Comentar para la preparacion de seminario
\include{C6Conclusiones} %Comentar para la preparacion de seminario
\include{C8Recomendación} %Comentar para la preparacion de seminario

%Generacion automatica de la lista de referencias 
\bibliographystyle{apacite} %estilo APA
%\bibliographystyle{unsrt} %estilo tipo EEE 
\bibliography{references} %indica nombre del archivo de referencias
\appendix
\include{C7Anexos} % En caso que sea necesario, comentar para la preparacion de seminario

\end{document}
